\documentclass[12pt]{article}
\usepackage{amsmath, amssymb}
\usepackage{array}
\usepackage{geometry}
\geometry{margin=1in}

\begin{document}

\title{CSE 400: Fundamentals of Probability in Computing}
\author{}
\date{}
\maketitle

\section*{Lecture Scribe — Joint Random Variables}

\hrule
\vspace{0.5cm}

\section*{Why Study Joint Random Variables}

A random variable assigns a numerical value to the outcome of a random experiment. In many situations, more than one quantity arising from the same experiment is random. In such cases, a single random variable does not fully describe the outcome of the experiment.

To represent multiple random quantities simultaneously, two or more random variables are defined and studied together. The study of these variables together is carried out using \textbf{joint random variables}. Joint random variables allow probabilities to be assigned to combinations of outcomes of multiple random variables.

\hrule
\vspace{0.5cm}

\section*{Discrete Case}

\subsection*{Joint PMF Table Representation}

Let $(X)$ and $(Y)$ be two discrete random variables. The \textbf{joint probability mass function (joint PMF)} of $(X)$ and $(Y)$ is defined as

\[
p_{X,Y}(x,y) = P(X = x, Y = y)
\]

The joint PMF gives the probability that the random variable $(X)$ takes the value $(x)$ and the random variable $(Y)$ takes the value $(y)$ simultaneously.

The joint PMF can be represented in a table form. In such a table:

\begin{itemize}
\item Rows correspond to possible values of $(X)$.
\item Columns correspond to possible values of $(Y)$.
\item Each entry represents the probability $(P(X=x, Y=y))$.
\end{itemize}

The probabilities must satisfy the normalization condition

\[
\sum_x \sum_y p_{X,Y}(x,y) = 1
\]

where the summations are taken over all possible values of $(X)$ and $(Y)$.

\hrule
\vspace{0.5cm}

\subsection*{Example: Rolling Two Dice}

Consider an experiment in which two fair six-sided dice are rolled.

Define the random variables

\[
X = \text{Outcome of the first die}
\]

\[
Y = \text{Outcome of the second die}
\]

Each die can take the values

\[
1,2,3,4,5,6
\]

The total number of ordered outcomes is

\[
6 \times 6 = 36
\]

Because the dice are fair, each ordered pair $((x,y))$ occurs with probability

\[
P(X=x, Y=y) = \frac{1}{36}
\]

\subsection*{Joint PMF Table}

\[
\begin{array}{c|cccccc}
(X \backslash Y) & 1 & 2 & 3 & 4 & 5 & 6 \\
\hline
1 & (1/36) & (1/36) & (1/36) & (1/36) & (1/36) & (1/36) \\
2 & (1/36) & (1/36) & (1/36) & (1/36) & (1/36) & (1/36) \\
3 & (1/36) & (1/36) & (1/36) & (1/36) & (1/36) & (1/36) \\
4 & (1/36) & (1/36) & (1/36) & (1/36) & (1/36) & (1/36) \\
5 & (1/36) & (1/36) & (1/36) & (1/36) & (1/36) & (1/36) \\
6 & (1/36) & (1/36) & (1/36) & (1/36) & (1/36) & (1/36) \\
\end{array}
\]

\subsection*{Verification of Total Probability}

The sum of all entries must equal 1.

\[
\sum_{x=1}^{6}\sum_{y=1}^{6} P(X=x,Y=y)
\]

There are 36 entries and each equals $(1/36)$.

\[
36 \times \frac{1}{36} = 1
\]

Thus the normalization condition is satisfied.

\hrule
\vspace{0.5cm}

\section*{Continuous Case}

For continuous random variables $(X)$ and $(Y)$, the joint distribution is described using a \textbf{joint probability density function (joint PDF)}.

The joint PDF is denoted as

\[
f_{X,Y}(x,y)
\]

The probability that the pair $((X,Y))$ lies within a region $(R)$ is given by

\[
P((X,Y) \in R) =
\int\int_R f_{X,Y}(x,y),dx,dy
\]

The joint PDF must satisfy

\[
\int_{-\infty}^{\infty} \int_{-\infty}^{\infty} f_{X,Y}(x,y),dx,dy = 1
\]

\subsection*{Geometric Interpretation}

In the continuous case, the joint PDF can be interpreted geometrically.

The function $(f_{X,Y}(x,y))$ forms a surface above the $(xy)$-plane. The probability that the random pair $((X,Y))$ lies in a region $(R)$ is equal to the \textbf{volume under the surface $(f_{X,Y}(x,y))$} above that region.

Thus,

\[
P((X,Y)\in R)
\]

corresponds to the volume under the surface above region $(R)$.

\subsection*{Worked Example}

Consider the joint probability density function

\[
f_{X,Y}(x,y)=
\begin{cases}
c & 0 \le x \le 1,; 0 \le y \le 1 \\
0 & \text{otherwise}
\end{cases}
\]

\subsubsection*{Step 1: Determine $(c)$}

The total probability must equal 1.

\[
\int_0^1 \int_0^1 c,dx,dy = 1
\]

First integrate with respect to $(x)$:

\[
\int_0^1 c,dx = c(1-0) = c
\]

Then integrate with respect to $(y)$:

\[
\int_0^1 c,dy = c(1-0) = c
\]

Thus,

\[
c = 1
\]

Therefore,

\[
f_{X,Y}(x,y)=
\begin{cases}
1 & 0 \le x \le 1,; 0 \le y \le 1 \\
0 & \text{otherwise}
\end{cases}
\]

\hrule
\vspace{0.5cm}

\section*{Joint Cumulative Distribution Function}

\subsection*{Definition of Joint CDF}

For random variables $(X)$ and $(Y)$, the \textbf{joint cumulative distribution function (joint CDF)} is defined as

\[
F_{X,Y}(x,y) = P(X \le x, Y \le y)
\]

The joint CDF gives the probability that the random variable $(X)$ is less than or equal to $(x)$ and the random variable $(Y)$ is less than or equal to $(y)$.

\subsection*{Continuous Case}

If $(X)$ and $(Y)$ are continuous random variables with joint PDF $(f_{X,Y}(x,y))$, the joint CDF is

\[
F_{X,Y}(x,y)
============

\int_{-\infty}^{x}
\int_{-\infty}^{y}
f_{X,Y}(u,v),dv,du
\]

This expression accumulates the probability density over the region

\[
u \le x, \quad v \le y
\]

\subsection*{Discrete Case}

If $(X)$ and $(Y)$ are discrete random variables with joint PMF $(p_{X,Y}(x,y))$, the joint CDF is

\[
F_{X,Y}(x,y)
============

\sum_{u \le x}
\sum_{v \le y}
p_{X,Y}(u,v)
\]

The sums include all values satisfying

\[
u \le x,\quad v \le y
\]

\subsection*{Properties of the Joint CDF}

\subsubsection*{Property 1: Range}

\[
0 \le F_{X,Y}(x,y) \le 1
\]

\subsubsection*{Property 2: Nondecreasing Property}

The joint CDF is nondecreasing in each argument.

If

\[
x_1 \le x_2
\]

then

\[
F_{X,Y}(x_1,y) \le F_{X,Y}(x_2,y)
\]

If

\[
y_1 \le y_2
\]

then

\[
F_{X,Y}(x,y_1) \le F_{X,Y}(x,y_2)
\]

\subsubsection*{Property 3: Limits}

\[
\lim_{x \to -\infty} F_{X,Y}(x,y) = 0
\]

\[
\lim_{y \to -\infty} F_{X,Y}(x,y) = 0
\]

\[
\lim_{x \to \infty, y \to \infty} F_{X,Y}(x,y) = 1
\]

\hrule
\vspace{0.5cm}

\section*{Comprehensive Example}

Consider the experiment of rolling two fair six-sided dice.

Define

\[
X = \text{Outcome of the first die}
\]

\[
Y = \text{Outcome of the second die}
\]

Each ordered pair $((x,y))$ occurs with probability

\[
P(X=x,Y=y)=\frac{1}{36}
\]

The joint cumulative distribution function is

\[
F_{X,Y}(x,y)=P(X \le x, Y \le y)
\]

\subsection*{Compute $(F_{X,Y}(2,3))$}

\[
F_{X,Y}(2,3)=P(X \le 2, Y \le 3)
\]

Possible values satisfying the conditions:

\[
X = 1,2
\]

\[
Y = 1,2,3
\]

The number of ordered pairs is

\[
2 \times 3 = 6
\]

Each pair has probability

\[
\frac{1}{36}
\]

Thus,

\[
F_{X,Y}(2,3)
============

6 \times \frac{1}{36}

\frac{6}{36}

\frac{1}{6}
\]

\section*{Result}

\[
F_{X,Y}(2,3) = \frac{1}{6}
\]

This value represents the probability that the outcome of the first die is less than or equal to 2 and the outcome of the second die is less than or equal to 3.

\end{document}